% 11_conclusion.tex — Conclusion
% ============================================================

\section{Conclusion}
\label{sec:conclusion}

We began with a paradox: computational imaging possesses solvers of
extraordinary power---deep unfolding networks that converge in seconds,
diffusion priors that hallucinate plausible textures, implicit
representations that compress entire scenes into coordinate
functions---yet real-world deployments routinely fail.  Reconstructions
that achieve 35+\,dB on simulated benchmarks collapse by 3--14\,dB under
realistic operator mismatch (\cref{tab:sensitivity}).  The gap is not a solver problem.
It is an infrastructure problem.

For more than a decade, the field has been optimizing the wrong layer.  It has
poured resources into the \emph{train}---faster architectures, cleverer
loss functions, larger training sets---while neglecting the
\emph{rail}: the calibrated forward operators, standardized evaluation
protocols, and diagnostic tools that determine whether a solver's
theoretical performance survives contact with reality.

This paper introduced PWM, the Physics World Model, as the rail for
computational imaging.  The contributions are concrete and measurable:

\begin{itemize}
  \item \textbf{64 modalities} formalized as composable
        \texttt{OperatorGraph} templates, each encoding the physics of
        measurement rather than approximating it.
  \item \textbf{89 OperatorGraph templates} that decompose forward models
        into calibratable, swappable subgraphs---making operator
        correction a modular operation rather than a monolithic
        reimplementation.
  \item \textbf{The four-scenario evaluation protocol}
        (Ideal / Assumed / Corrected / Oracle~Mask), which for the first
        time separates solver error from operator error and quantifies
        the recovery ratio of any pipeline.
  \item \textbf{LIP-Arena and the Triad Law}, providing real-time
        leaderboard ranking and a diagnostic framework that decomposes
        reconstruction failure into recoverability, carrier budget, and
        operator mismatch components.
\end{itemize}

The evidence is unambiguous.  Across the 26-modality benchmark,
operator correction alone improved reconstruction quality by $+0.54$ to
$+48.25$\,dB across 9 correction configurations spanning 7 distinct modalities (16 registered)---gains that no solver upgrade can
replicate.  The CASSI deep-dive was particularly revealing: the mismatch
penalty ($13.98$~dB) exceeded the solver upgrade gain ($10.47$~dB) by
$1.34\times$.  The bottleneck was never the
algorithm.  It was the physics encoding.

These results carry a temporal urgency.  The foundry window described in
\cref{sec:roadmap} is open now but will not remain so.  Within
18~months, the evaluation norms of computational imaging will
solidify---either around principled, physics-aware infrastructure, or
around the fragmented, simulation-only benchmarks that have defined the
field to date.  The QWERTY effect guarantees that whichever protocol is
adopted first by a critical mass of laboratories will become permanent.

\medskip
\noindent
A problem is solved when the bottleneck shifts from genius to compute.
PWM provides the infrastructure for that shift.
